\section{The Diagnosis}

\textbf{[K]} From A085 (the constants of nature $c, \hbar, k_B$) and A130
($\alpha$), a single sharp statement can be drawn that until now appeared only
scattered:
\begin{quote}
\emph{Every constant that can be set to one through a consistent unit or
normalisation choice is thereby proved to be \emph{not fundamental}. It is a
conversion factor, not a degree of freedom.}
\end{quote}

The settability to one is therefore not merely a convenient notation but a
\emph{test}. What can be brought to one without losing any physical statement was
never content of nature but a property of the measurement system used. The famous
numerical value -- $c = 299\,792\,458$, $1/\alpha = 137{.}036$ -- then measures
the \emph{historical unit}, not nature.

\section{Two Paths to One, One Result}

\textbf{[K]} There are two different reasons why a constant can be brought to one,
and it is important to distinguish them -- because the result is the same in both
cases, but the reason is not.

\textbf{[B]} \emph{First, via a physical identity.} $c$, $\hbar$, $k_B$ become
one because the time-mass duality \emph{identifies} their dimensions (A085): mass,
energy, temperature, and inverse time are one quantity in four guises. The $=1$ is
here the reading of an identity. The reason is substantive.

\textbf{[K]} \emph{Second, via a normalisation.} $\alpha$ becomes one because one
normalises the charge so that $e^2 = 4\pi$ (A130). The reason is a unit choice,
not a physical statement.

\textbf{[K]} The crucial point: \emph{for the diagnosis the distinction is
irrelevant}. Whether a constant falls to one via an identity or via a
normalisation -- in both cases the settability shows that it is not an independent
degree of freedom. The reason decides \emph{how deeply} the bookkeeping is anchored
(for the duality, in a physical identity; for $\alpha$, in a convention), not
\emph{whether} the constant is fundamental. Neither is fundamental.

\section{What Cannot Be Set to One}

\textbf{[K]} The statement is sharp only if its converse holds: there must be
quantities that \emph{no} unit choice can bring to one -- and precisely those are
the content.

\textbf{[B]} Two kinds survive the test:
\begin{itemize}
  \item \textbf{The geometric parameter $\xi$.} It is a pure numerical ratio of
    the packing geometry (A020). No change of the metre, the second, or the charge
    unit changes it; it is $4/30000$ in every system of units. $\xi$ cannot be
    normalised to one without changing the geometry itself.
  \item \textbf{Dimensionless ratios.} $m_\mu/m_e = 206{.}77$ is the quotient of
    two quantities of \emph{the same} dimension; no unit choice cancels it. Such
    ratios, unlike $\alpha$, are not reachable by a normalisation -- $\alpha$
    conceals a charge unit, a mass ratio conceals nothing.
\end{itemize}

\textbf{[STIPULATION]} Thus the dividing line is: settable to one $\Rightarrow$
convention; not settable $\Rightarrow$ content. $\alpha$ stands, despite its
dimensionlessness, on the convention side, because it carries a charge unit;
$m_\mu/m_e$ stands on the content side. The only remaining content of the theory
is $\xi$ and a single absolute scale.

\section{Two Sides of a Coin -- Yet Only One Fundamental}

\textbf{[K]} Here a contradiction threatens that must be cleared up. A130 calls
$\alpha$ and $\xi$ ``two sides of a coin'', a single degree of freedom, connected by
\[
  \alpha_\text{SI} = \xi\cdot E_0^2 .
\]
This document simultaneously says $\alpha$ is not fundamental, but $\xi$ is.
How do both fit together?

\textbf{[K]} They fit as soon as one reads ``one degree of freedom'' correctly:
the statement concerns the \emph{information content}, not the \emph{rank}.
$\alpha_\text{SI}$ and $\xi$ carry \emph{the same} geometric information -- one
who knows $\xi$ and the scale $E_0$ knows $\alpha_\text{SI}$, and vice versa. In
this sense it is one quantity, not two. That is the coin.

\textbf{[B]} But the two sides are not equally ranked. $\xi$ is this information
\emph{bare}: a pure ratio, $4/30000$ in every system of units. $\alpha_\text{SI}$
is \emph{the same} information \emph{dressed}: $\alpha_\text{SI} = \xi\cdot E_0^2$,
where the factor $E_0^2$ is pure scale -- the unit dress. The famous value $1/137$
is therefore not $\xi$ beside $\alpha$ but $\xi$ \emph{in the costume} of $E_0^2$.

\textbf{[K]} The diagnosis of this document confirms exactly that. Strip the units
-- set $\alpha = 1$ in natural units -- and $\alpha$ carries nothing more: the
information has not disappeared, it sits entirely in $\xi$ and the one scale.
$\xi$ survives every unit change; $\alpha = 1$ is empty afterwards. The coin
therefore remains intact, but it \emph{is} $\xi$: $\alpha_\text{SI}$ is how this
coin looks in SI light, not a second coin.

\textbf{[STIPULATION]} In short: ``two sides'' means same content, not same rank.
One degree of freedom -- and its fundamental form is $\xi$, its dressed form
$\alpha_\text{SI}$. The apparent contradiction arises only when one confuses
``equivalent'' with ``equally fundamental''.

\section{Why SI Computation at All?}

\textbf{[K]} If in natural units everything except $\xi$ and one scale falls to
one, the question arises why one needs the more elaborate SI calculations, the
bridge constants, and the multiple derivation paths at all. The answer separates
three things that are easily confused.

\textbf{[K]} \emph{First: translation.} Measurement happens in the laboratory, and
laboratories compute in SI. The SI dress is the \emph{translation} of the geometric
core into readable numbers -- indispensable, but not itself physics. One needs it
because the measuring device speaks in metres and seconds, not because the content
demands it. A result in natural units is complete; the SI number is its
pronunciation in a particular language.

\textbf{[B]} \emph{Second: overdetermination as test.} The multiple derivation
paths add nothing to the content -- they \emph{test} it. If two paths lead to the
same energy scale and their difference is exactly the derived fractal factor (A130),
or if three derivations of the lepton masses hit the same value (A110), that is not
an additional degree of freedom but a cross-check. A single path could be a
disguised fit; several independent paths that coincide cannot be. The redundancy is
the hardening, not the content.

\textbf{[K]} \emph{Third: connectivity.} Different paths bind the one $\xi$-content
to different known physics -- Yukawa coupling, Koide relation, RG running. This
makes the same core testable from several familiar directions and prevents it from
appearing consistent only in its own language. The content remains one; the
approaches are many so that it can be controlled from outside.

\textbf{[STIPULATION]} In short: the SI calculation is translation, the multiple
paths are control. Both are needed to connect the geometric core to measurement
and to harden it -- neither adds a degree of freedom to the core. The core remains
$\xi$ and one scale.

\section{Verification Script}

\begin{itemize}
  \item The diagnosis automated: for $c, \hbar, k_B, \alpha$ it is shown that a
    unit/normalisation choice brings them to one (hence convention); for $\xi$ and
    $m_\mu/m_e$, that no unit choice achieves this (hence content); and the
    cross-check that two paths hit the same value without adding a degree of
    freedom:
    \texttt{python/A\_Serie\_Skripte/a135\_eins\_diagnose.py}
\end{itemize}

\section{Open Questions}

\textbf{First, the one scale remains an input.} The statement reduces everything to
$\xi$ and one absolute scale; what value this scale has is itself not derived (A130)
but the one remaining calibration point. The scale is \emph{not} settable to one
in the sense that its value carries content -- but its numerical value is the one
input the system needs.

\textbf{Second, the diagnosis is not a proof of completeness.} That everything
\emph{so far} either falls to one or is $\xi$ shows the parsimony of the theory;
it does not prove that no further independent quantity will ever be needed. The
statement orders what is present; it does not exclude what is future.

\section{Supersedes}

This document subsumes: Dok.~043 and Dok.~044 (dissolution of constants), Dok.~101
(circularity of constants, one degree of freedom). It draws together the diagnosis
carried out case by case in A085 (duality) and A130 ($\alpha$), and answers the
question of the purpose of SI computation. Incorporated: P40.

\section{Use of AI Tools}

This document was created with the assistance of AI tools. The questions,
stipulations, and all substantive decisions come from the author; the tools
formulate, compute, and create verification scripts.
Responsibility for content and errors rests entirely with the author.
Full wording of the rules in A010.
